Language models can learn new relational structure from context, but a learned internal geometry is useful to a user only if it can be addressed through ordinary prompts. We test this behavioral question for graph geometries inspired by recent work on in-context representation learning. We assign ordinary words to hidden nodes of either a 4x4 grid or a 16-node ring, expose the model to random-walk traces over the graph, and ask \gptfourone to answer bundled word-adjacency questions. We vary graph type, trace length, and prompt style, comparing direct JSON-only answers against an explicit-reasoning ablation and against an observed-edge memorization baseline. Across 1,280 scored query-level examples from 160 API calls, direct prompting reaches 72.7\% accuracy and explicit reasoning reaches 73.9\%, both below the 82.5\% observed-edge baseline. The critical failure is sparse topology: on true adjacencies never observed in the prompt, direct prompting answers correctly in only 4/112 cases (3.6\%), while explicit reasoning answers 0/112 correctly. The results show that these repeated novel geometries are verbally addressable mainly as memories of transitions already present in context. They do not support the stronger claim that ordinary prompting exposes unobserved parts of a newly induced topology under limited repetition.
